% Chapter 5: winding step by step.
\chapter{Winding paper step by step}\label{ch:winding}
\index{winding}

\section{Before you start}
\begin{itemize}
  \item Power off. Load the paper supply and thread the paper under the
        measuring wheel and onto the take-up core, as you were shown at
        installation ([THREADING DIAGRAM -- SEE INSTALLER NOTES]).
  \item Check that the measuring wheel sits on the paper and turns freely.
  \item Power on. Wait about one second. The idle screen appears and
        \led{1} and \led{3} light.
\end{itemize}

\begin{lcd}
 -= PAPER WINDER =-

[1] 30yd    [2] 90yd
[4] click = manual
\end{lcd}

\section{Start a wind}
Push \btn{1} for a \textbf{30 yard} roll or \btn{2} for a \textbf{90 yard}
roll. One push is enough. Starting a wind also silences any buzzer that was
still sounding.

The relay closes, the motor runs and the screen shows the progress:

\begin{lcd}
     WINDING...
 45% [#####--------]
LEN   12.3/  30.0 yd
[1/2]STOP   [3]PAUSE
\end{lcd}

The LEDs fill up in quarters as the wind goes on. If the machine has seen
encoder errors, the top row shows \shows{WINDING    ERR:n} instead (see
Section~\ref{sec:screen-winding}).

\section{The end of a wind: coast and trim}\label{sec:trim}
\index{trim}\index{coast}
The motor cannot stop the roll instantly. After the relay opens the roll
keeps turning for a short distance, called the \textbf{coast} (set at the
factory to 12 inches). So the machine switches the motor off about one coast
\emph{before} the target, and lets the roll glide up to it.

Then it \textbf{trims}:
\begin{enumerate}
  \item The screen shows \shows{TRIMMING...}. The machine waits half a second
        for the paper to stop.
  \item If the length is not yet enough, it gives the motor a short bump
        and waits again. It repeats this until the length is at least the
        target plus the safety margin (1 inch from the factory).
  \item When the length is enough, it beeps \textbf{three short beeps} and
        shows \shows{DONE}.
\end{enumerate}

\begin{lcd}
     TRIMMING...
 99% [############-]
LEN   29.9/  30.0 yd
[1/2]STOP   [3]PAUSE
\end{lcd}

\section{Finished}
\begin{lcd}
 -= PAPER WINDER =-
DONE 30.1/30.0 yd
+4.3 in over
[1] 30yd    [2] 90yd
\end{lcd}

The second row shows the measured length and the target. The third row says
how far over (or short) the roll is compared with the target, in inches (or in
yards when it is a yard or more). Just after the motor stops the screen shows
\shows{NOW} and \shows{settling...} for a moment while the roll comes to rest;
then \shows{DONE} and the final figure. Cut the paper and remove the roll. The result stays on the screen
for 60 seconds, then the normal idle screen returns. You can start the next
wind straight away; a button still sitting pushed in is fine.

\section{Pause, nudge and resume}\label{sec:pause}
\index{pause}\index{nudge}
\begin{itemize}
  \item Push \btn{3} to \textbf{pause}. The motor stops and the screen shows
        \shows{PAUSED}. The bottom row changes to \shows{[3]PLAY}.
  \item While paused, push \btn{4} to run the motor by hand (the screen shows
        \shows{PAUSED - MANUAL RUN}); push \btn{4} again to stop. The motor
        also stops by itself if the length reaches the target.
  \item Push \btn{3} again to \textbf{resume}. The wind carries on from where
        it was; the hand-run length counts towards the roll.
\end{itemize}

\section{Stop a wind early}\label{sec:stop}
\index{stop}
Push \btn{1} or \btn{2}. The motor stops at once. What you see next:
\begin{itemize}
  \item If the roll had already reached the full length plus margin, it is
        reported as \shows{DONE} as usual.
  \item Otherwise the screen shows \shows{!! UNDER TARGET !!} with
        \shows{STOPPED} and how much is missing. There is no alarm sound for a
        stop you made yourself. The screen stays until you push a button.
\end{itemize}

\section{When the roll comes up short}\label{sec:under-target}
\index{under target}\index{short roll}
If the trim cannot finish (the paper does not move when bumped, too many
bumps, or the trim takes longer than 30 seconds), the machine gives up and
sounds the \textbf{under-target alarm}: three long beeps, repeated every 6
seconds for about a minute.

\begin{lcd}
 !! UNDER TARGET !!
SHORT -8.2 in
LEN 29.8/30.0 yd
[1] 30yd    [2] 90yd
\end{lcd}

What to do:
\begin{enumerate}
  \item Push any button to silence the alarm. (That push also does its
        normal job, so \btn{1} would start a new 30 yard wind. Push \btn{3}
        if you only want to clear the screen.)
  \item Switch the power off and look for the cause: paper run out, a jam,
        the wheel lifted off the paper, a torn web.
  \item Do \textbf{not} sell the short roll as full length. Either rewind it
        or add the missing length by hand.
  \item See Chapter~\ref{ch:trouble} if it keeps happening.
\end{enumerate}

\section{Manual mode}\label{sec:manual}
\index{manual mode}
From the idle screen, push \btn{4}: the motor runs continuously and the
screen shows \shows{MANUAL: motor ON}; only \led{4} is lit. Push \btn{4}
again to stop. Use this for threading and for winding odd lengths by eye. The
machine does not count or stop a manual run for you.
\begin{warning}
Never leave the machine in manual mode unattended.
\end{warning}

\section{After a power cut or reset}
If the power fails, or the controller resets itself (watchdog), the relay
opens, the motor stops, and the machine starts up at the idle screen. Any
wind in progress and its result are lost. Measure the roll by hand or rewind
it.
